% NSDI 2027 Paper Template
% The 24th USENIX Symposium on Networked Systems Design and Implementation
% May 11-13, 2027, Providence, RI, USA
%
% Format: ≤12 pages (excluding references), USENIX format
%         Two-column, 10pt on 12pt leading, Times Roman
%
% Official CFP: https://www.usenix.org/conference/nsdi27/call-for-papers
% Template source: https://www.usenix.org/conferences/author-resources/paper-templates
%
% IMPORTANT NOTES:
% - Three tracks: Traditional Research, Frontiers, Operational Systems
% - Indicate track on title page and submission form
% - Prescreening phase: reviewers read ONLY the Introduction
% - Introduction must articulate track-specific criteria
% - Two deadlines: Spring (April 2026) and Fall (September 2026)
% - One-shot revision available for rejected papers
%
% TRACK REQUIREMENTS IN INTRODUCTION:
% - Research: articulate novel idea + evaluation evidence
% - Frontiers: articulate novel non-incremental idea
% - Operational: describe deployment setting, scale, lessons learned

\documentclass[letterpaper,twocolumn,10pt]{article}
\usepackage{usenix-2020-09}

% Recommended packages
\usepackage[utf8]{inputenc}
\usepackage{amsmath,amssymb}
\usepackage{graphicx}
\usepackage{booktabs}
\usepackage{hyperref}
\usepackage{url}
\usepackage{xspace}

% Custom commands
\newcommand{\system}{SystemName\xspace}  % Replace with your anonymized system name
\newcommand{\eg}{e.g.,\xspace}
\newcommand{\ie}{i.e.,\xspace}
\newcommand{\etal}{\textit{et al.}\xspace}

\begin{document}

% Indicate your track in the title page
% Options: [Research Track] / [Frontiers Track] / [Operational Systems Track]
\title{Your Paper Title Here}

\author{Paper \#XXX}  % Anonymized for submission (double-blind)
% Operational Systems track: may keep real company/system names for context
% Camera-ready:
% \author{
%   Author One\\
%   \textit{Affiliation One}\\
%   \texttt{email@example.com}
%   \and
%   Author Two\\
%   \textit{Affiliation Two}\\
%   \texttt{email@example.com}
% }

\maketitle

\begin{abstract}
% Write your abstract here.
Your abstract goes here.
\end{abstract}

\section{Introduction}
\label{sec:intro}

% CRITICAL: This section is used for prescreening!
% Reviewers will read ONLY this section to determine:
% 1. Subject falls within NSDI scope (networked/distributed systems)
% 2. Exposition understandable by NSDI PC member
% 3. Track-specific criteria met:
%    - Research: novel idea + evaluation evidence
%    - Frontiers: novel non-incremental idea
%    - Operational: deployment setting, scale, lessons
%
% Make sure your Introduction clearly articulates all of the above.

Your introduction goes here. Clearly state the problem, your approach,
and your contributions. Remember: this section is used for prescreening.

\section{Background and Motivation}
\label{sec:background}

Provide necessary background on networked systems context.

\section{Design}
\label{sec:design}

Describe the design of your networked system.

\section{Implementation}
\label{sec:implementation}

Describe implementation details.

\section{Evaluation}
\label{sec:evaluation}

Present your evaluation results.

\subsection{Experimental Setup}

Describe your experimental setup, workloads, and baselines.

\subsection{Results}

Present and discuss your results.

\section{Related Work}
\label{sec:related}

Discuss related work in networked systems.

\section{Conclusion}
\label{sec:conclusion}

Summarize your contributions and key findings.

% Bibliography
{\footnotesize \bibliographystyle{acm}
\bibliography{references}}

\end{document}
